\section{Related Work}

\paragraph{Tsunami and shallow-water numerical solvers.}
Tsunami propagation is commonly modeled with depth-averaged shallow-water equations, where finite-volume methods and Riemann-solver-based fluxes provide a robust framework for wave propagation over nonuniform bathymetry \citep{leveque2002finite,toro2009riemann}. Operational and research tsunami codes such as GeoClaw demonstrate the importance of well-balanced finite-volume methods, wet--dry handling, and adaptive refinement for transoceanic propagation and coastal inundation \citep{leveque2011tsunami,berger2011geoclaw}. In our benchmark, the primary reference solver follows the hydrostatic-reconstruction idea of \citet{audusse2004hydrostatic}, which preserves the lake-at-rest equilibrium and nonnegative water depth. We additionally include a second-order MUSCL-HR reference motivated by finite-volume tsunami simulation with hydrostatic reconstruction \citep{clain2016secondorder}, while using Rusanov-type numerical fluxes and wave-speed estimates in the spirit of standard hyperbolic finite-volume methods \citep{toro2009riemann,toro2026rusanov}. We also include a Boussinesq-type weakly dispersive comparison reference, since Boussinesq-type solvers provide context for modeling dispersive wave effects beyond non-dispersive shallow-water theory \citep{shi2012funwave,charlier2022boussinesq}.

\paragraph{Synthetic bathymetry and source generation.}
Because the goal of this work is a controlled benchmark rather than a site-specific operational model, we procedurally generate bathymetry and initial disturbance fields. Procedural terrain modeling and fractal/noise-based geometry generation have long been used to create diverse spatial fields \citep{musgrave1989terrain}. We use these ideas only as a controllable source of synthetic bathymetric variability, not as a replacement for real bathymetric products. The source generator includes simple Gaussian, dipole, rough, and Okada-like families; the latter is motivated by classical elastic half-space deformation models for earthquake sources \citep{okada1985surface}. Real global marine gravity and bathymetry products \citep{sandwell2014gravity} are relevant to future site-specific extensions, but this paper intentionally focuses on synthetic scenarios for controlled generalization tests.

\paragraph{Neural operators and baseline architectures.}
Neural operators learn maps between function spaces and are now standard tools for PDE surrogate modeling \citep{li2021fno,kovachki2023neuraloperator}. We use FNO as the main neural-operator surrogate and compare it with simpler convolutional baselines such as CNNs and U-Nets \citep{lecun1998gradient,ronneberger2015unet}. Later operator variants include factorized Fourier operators, U-FNO-style multiscale Fourier/U-Net hybrids, and wavelet neural operators, which motivate our baseline coverage rather than a broad architecture survey \citep{tran2023ffno,wen2022ufno,tripura2023wno}. Our contribution is not a new architecture; instead, we evaluate established surrogate families under a controlled multi-fidelity tsunami-like benchmark. Recent neural-operator work on localized kernels and geophysical autoregressive forecasting further motivates evaluating operator models under local structure, multi-step forecast horizons, and resolution changes \citep{liu2024localized,pathak2022fourcastnet,hu2025spherical}.

\paragraph{Tsunami-oriented neural surrogates.}
Machine-learning approaches have recently been explored for tsunami propagation, inundation, and rapid assessment. Closest to our setting, \citet{kim2026ffno_tsunami} propose a Factorized Fourier Neural Operator surrogate for basin-scale tsunami propagation using COMCOT-generated wavefields, held-out source scenarios, arrival-time errors, and speed comparisons. PINN and DeepONet-style models have also been studied for tsunami inundation modeling \citep{brecht2025pinn_tsunami}, while neural-network surrogates have been used to accelerate probabilistic tsunami inundation assessment \citep{fukutani2025prob_tsunami}. Recent work has also considered neural operators for rapid tsunami propagation and source inversion \citep{someya2026tsunamiNO}. These studies motivate neural surrogates for tsunami-related modeling, but our contribution is different: we construct a controlled benchmark with synthetic bathymetry/source generators, multiple shallow-water numerical references, explicit OOD and cross-resolution protocols, solver-vs-solver analysis, and uncertainty evaluation.

\paragraph{Scientific machine-learning benchmarks.}
Our work is also related to benchmark-oriented scientific machine learning. PDEBench provides a broad benchmark suite for time-dependent PDE problems, including generated datasets, baseline models, and multiple evaluation metrics \citep{takamoto2022pdebench}. In contrast, our benchmark is narrower but domain-focused: we target tsunami-like shallow-water propagation with controlled bathymetry/source variability and multi-solver reference labels.

\paragraph{Solver-fidelity and emulator-superiority-style evaluation.}
A central difficulty in PDE surrogate benchmarking is that the training solver is not necessarily ground truth. Recent work on neural emulator superiority shows that an emulator trained on lower-fidelity numerical data can, in specific regimes, be closer to a higher-fidelity reference than the solver that generated its training data \citep{koehler2026emulator}. This motivates our multi-reference evaluation: we report same-solver surrogate error, raw solver-vs-solver gaps, cross-solver surrogate evaluation, and emulator-superiority-style ratios only when the normalization and reference choices are physically meaningful. This separates surrogate approximation error from differences among numerical reference solvers.

\paragraph{Uncertainty quantification.}
For hazard-oriented surrogate models, point accuracy alone is insufficient; users also need to know when the model may be unreliable. Deep ensembles provide a simple and scalable approach to predictive uncertainty estimation and often give strong calibration without major changes to the base model \citep{lakshminarayanan2017deepensembles}. We therefore use ensembles as the primary uncertainty mechanism and evaluate whether uncertainty correlates with actual error and increases under OOD scenarios. Broader surveys and Bayesian deep-learning perspectives provide context for interpreting epistemic and aleatoric uncertainty in deep models \citep{he2023uqsurvey,wilson2020bayesian}.

\paragraph{Inverse problems.}
Inverse source reconstruction is an adjacent direction, with classical adjoint inversion and recent neural inverse-operator methods providing relevant foundations. We leave this task to follow-up work and focus here on forward emulation and evaluation.
